\section{A.4. Strict boundedness, fibre ranks, and perfection}
\label{sec:A4}

Let $L$ be a simplicial $R$-module and write $N(L)$ for its
nonnegative normalized chain complex.  ``$\partial$-perfect'' means
that $N(L)$ is quasi-isomorphic to a bounded complex of finite
projective $R$-modules.  It does not mean that the actual normalized
terms are bounded.  ``Strictly $\partial$-perfect'' does mean that
these actual terms are finite projective and zero outside a finite
interval.  This is the distinction made in 137:14, Definition~1.

\paragraph{A4.1 (Known, with its form here derived).}
\label{A4.1}
For every commutative $R$ and simplicial $R$-module $L$, there is a
natural decomposition
\[
 L_n\cong\bigoplus_{i=0}^{n}N_i(L)^{\oplus\binom ni}.
 \tag{A4.1a}
\]
In particular $L_n$ is finite projective for every $n$ if and only if
every $N_i(L)$ is finite projective; the same equivalence holds for
finite generation.  This is Dold--Puppe, equation~(3.3),
Theorem~3.6 and Corollary~3.8, pp.~219--220
\cite{DoldPuppe1961}.

\begin{proof}[Derivation of the stated form]
Equation~(3.3) of the opened source gives one copy of $N_i(L)$ for
each increasing surjection $[n]\twoheadrightarrow[i]$, and
Theorem~3.6 identifies its denormalization naturally with $L$.
Such a surjection is specified by its $i$ jump positions among the
$n$ gaps between successive elements of $[n]$, so there are
$\binom ni$ of them.  This proves (A4.1a).  Its finite direct sums
and its summand with $i=n$ prove both finiteness assertions.  This
also checks the base-change used below: tensoring (A4.1a) with any
$R$-algebra preserves the direct-sum decomposition, without a
flatness assumption.
\end{proof}

\paragraph{A4.2 (Proved; stronger than the regular-noetherian correction).}
\label{A4.2}
Suppose every $L_n$ is finitely generated over an arbitrary
commutative ring $R$.  For an integer $D\geq0$, the following are
equivalent:
\begin{enumerate}
\item $N_i(L)=0$ for every $i>D$.
\item For every prime ideal $\mathfrak p$, the function
\[
 f_{\mathfrak p}(n)=
 \dim_{\kappa(\mathfrak p)}(L_n\otimes_R\kappa(\mathfrak p))
\]
agrees, for all sufficiently large $n$, with a polynomial of degree
at most $D$.  The starting value of $n$ may depend on $\mathfrak p$.
\end{enumerate}
Maximal ideals alone suffice in the second condition.  When these
conditions hold the polynomial identity is valid for every
$n\geq0$, with
\[
 f_{\mathfrak p}(n)=\sum_{i=0}^{D}
 a_i(\mathfrak p)\binom ni,
 \qquad a_i(\mathfrak p)=
 \dim_{\kappa(\mathfrak p)}(N_i(L)\otimes_R\kappa(\mathfrak p)).
 \tag{A4.2a}
\]
For a nonzero fibre its degree is the largest $i$ with
$a_i(\mathfrak p)\ne0$.  The zero fibre has the zero polynomial.
Regularity and noetherianness are not needed for this equivalence.

\begin{proof}
Finite generation of $N_i(L)$ follows from A4.1.  Tensoring the
direct-sum formula with a residue field gives
$f_{\mathfrak p}(n)=\sum_{i=0}^{n}\binom ni a_i(\mathfrak p)$,
with nonnegative integer coefficients.  Bounded normalization
therefore implies (A4.2a).  Conversely suppose the eventual degree
is at most $D$.  If $a_j(\mathfrak p)>0$ for some $j>D$, then for
every $n\geq j$ we have
$f_{\mathfrak p}(n)\geq\binom nj$.  This contradicts its eventual
polynomial growth of degree at most $D$.  Thus
$N_j(L)\otimes_R\kappa(\mathfrak p)=0$ for every $j>D$ and every
prime (or maximal) ideal.  Local Nakayama applied to each finitely
generated $N_j(L)$ makes all its maximal localizations zero;
hence $N_j(L)=0$.  Finally the leading coefficient of a nonzero
sum (A4.2a), with largest index $j$, is
$a_j(\mathfrak p)/j!>0$, proving the assertion about the degree.
All arguments include $D=0$.  For the zero ring all modules and
all relevant functions are zero and the equivalence still holds.
\end{proof}

\paragraph{A4.3 (Known, with the requested corrected consequence proved).}
\label{A4.3}
Let $R$ be regular noetherian, with no finite-Krull-dimension
assumption.  For any simplicial $R$-module $L$,
\[
 L\text{ is }\partial\text{-perfect}
 \quad\Longleftrightarrow\quad
 \pi_i(L)\text{ is finite for every }i
 \text{ and }\pi_i(L)=0\text{ for }i\gg0.
 \tag{A4.3a}
\]
This is the normalized-complex form of the Stacks Project,
Lemma~15.76.14, Tag~066Z \cite{Stacks066Z}.  Consequently, under
the termwise finiteness assumption in A4.2, the uniformly bounded
fibre-polynomial condition implies $\partial$-perfection.  It is
equivalent to actual bounded normalization, not to perfection.

\begin{proof}[Derivation and verification of the hypotheses]
The Dold--Puppe equivalence identifies $\pi_i(L)$ with
$H_i(N(L))$.  Reindexing cohomologically by
$K^{-i}=N_i(L)$ converts bounded homology and finite homology
modules exactly into the conditions of Tag~066Z.  Its conclusion
is exactly the definition of $\partial$-perfection.  To apply it
to A4.2, bounded normalization and noetherianness imply finite,
bounded homology, so all its hypotheses hold.

We also verify the potentially delicate dimension issue.  The
opened proof of Tag~066Z does not impose a uniform bound on the
dimensions of all local rings.  For a fixed finite module it
takes a resolution by finite free modules, lets $U_j$ be the open
set where its $j$th syzygy is locally free, and uses
$U_j\subseteq U_{j+1}$.  Regularity of every localization makes
these opens cover the spectrum; quasi-compactness gives a finite
stage covering the spectrum.  Thus that particular finite module
has a finite projective resolution.  A bounded complex with
finite homology then has a finite filtration by shifts of such
modules in the derived category.  This explains why arbitrary
regular noetherian rings, including those of unbounded Krull
dimension, are within the claim.
\end{proof}

\paragraph{A4.4 (Refuted; both directions of 137:14--15, Proposition~1(e)).}
\label{A4.4}
The original equivalence between $\partial$-perfection and
termwise finiteness plus eventual polynomial fibre ranks is false
in both directions, even for $R=\mathbf Z$.  The added uniform
degree bound repairs the implication to perfection, but it cannot
repair the converse.  The strongest direct rank statement is A4.2;
the derived statement is A4.3.

\begin{proof}
For failure of sufficiency, let $p_1,p_2,\ldots$ be distinct
rational primes and let the nonnegative chain complex $M$ have
$M_i=\mathbf Z/p_i$ for $i\geq1$, $M_0=0$, and all differentials
zero.  Let $L$ be its Dold--Puppe denormalization.  Each $L_n$ is a
finite direct sum of finite groups, hence is finitely generated.
At the prime $(p_j)$ its fibre dimension is $\binom nj$, for all
$n$, because every $\mathbf Z/p_i$ with $i\ne j$ tensors to zero.
At the generic point it is zero.  Thus every fibre-rank function
is polynomial, but the degrees are unbounded as the prime varies.
The homology $H_i(M)=\mathbf Z/p_i$ is nonzero in every positive
degree, so $M$ cannot be quasi-isomorphic to a bounded projective
complex.  This proves nonperfection.

For failure of necessity let $M=\bigoplus_{i\geq0}D_i$, where the
sum is locally finite in chain degree.  In degree zero its rank
is one; in each positive degree its rank is two.  Each summand
has its evident contraction, and their locally finite direct sum
contracts $M$.  Hence $M$ is quasi-isomorphic to zero and is
perfect.  Its denormalization $L$ is termwise finite free, but
A4.1 and the binomial theorem give
\[
 \operatorname{rank}L_n
 =1+2\sum_{i=1}^{n}\binom ni
 =2^{n+1}-1.
\]
This is not eventually polynomial.  For example, for every fixed
$D$ the ratio $2^n/n^D$ tends to infinity, as follows by taking
the ratio of successive terms or by the binomial lower bound of
degree $D+1$.  The counterexample applies over every nonzero
commutative base ring.  The attached program checks the exact
rank identities in degrees $0$ through $16$; the argument above,
not the finite check, proves them in all degrees.
\end{proof}

\paragraph{A4.5 (Known criterion, and nonregular sufficiency refuted).}
\label{A4.5}
Let $R$ be any noetherian ring.  For a simplicial $R$-module $L$
whose normalized complex is pseudo-coherent, perfection is
equivalent to the existence of a single $D\geq0$ such that
\[
 H_i\bigl(N(L)\otimes_R^{\mathbf L}\kappa(\mathfrak p)\bigr)=0
       \quad\text{for every }i>D
       \text{ and every prime }\mathfrak p.
 \tag{A4.5a}
\]
Maximal ideals suffice.  A sufficient hypothesis for the required
pseudo-coherence is that all $L_n$ be finite; more generally it
suffices that all $\pi_i(L)$ be finite.  Formula (A4.5a) is the
specialization of the Stacks Project, Lemma~15.79.2,
Tag~0G99 \cite{Stacks0G99}, to a nonnegative chain complex.
Thus on a nonregular ring A4.2 remains the exact answer about
actual boundedness; derived residue-field fibres give the exact
additional test for perfection.  Bounded normalization alone
does not imply perfection.

\begin{proof}[Derivation, and complete counterexample]
Over a noetherian ring, a nonnegative chain complex with finite
homology in every degree has a quasi-isomorphism from a
nonnegative degreewise finite-free complex.  One constructs it
inductively: choose finitely many cycles generating the lowest
homology, map a finite free module onto them, then choose
finitely many generators of the next homology and of the kernel
created at the preceding step.  The needed kernels are finite
by noetherianness, and this construction successively makes the
cone acyclic in degrees $0,1,\ldots$.  This is pseudo-coherence
after cohomological reindexing.  Termwise finite modules give
finite homology by noetherianness, so they imply the same
hypothesis.

Put $K^{-i}=N_i(L)$.  Its derived tensor with any residue field
has no positive cohomology, because it can be computed by the
nonnegative free chain resolution just described.  Thus
(A4.5a) says exactly that its derived fibres have cohomology only
in $[-D,0]$.  Lemma~15.79.2 applies to this pseudo-coherent $K$
and states that this is equivalent to perfection with Tor
amplitude in that interval.  Conversely a perfect nonnegative
complex has some finite Tor bound and therefore satisfies the
same condition.  This explicitly checks the uniformity and the
derived, rather than underived, tensor product in the cited
criterion.

For a finite counterexample to boundedness implying perfection,
take $R=\mathbf F_2[\epsilon]/(\epsilon^2)$ and
$M=R/(\epsilon)=\mathbf F_2$, and let $L$ be the constant
simplicial module $M$.  Its normalization is $M$ in degree zero
and zero elsewhere, and its sole fibre-rank function is the
constant polynomial one.  The ring is noetherian and nonregular:
it is local of Krull dimension zero with nonzero maximal ideal,
so it is not a field.  Its infinite free resolution is
\[
 \cdots\xrightarrow{\epsilon}R
       \xrightarrow{\epsilon}R
       \xrightarrow{\epsilon}R\longrightarrow M\longrightarrow0.
\]
For $a+b\epsilon\in R$ multiplication by $\epsilon$ is
$a\epsilon$, so its image and kernel are both $(\epsilon)$;
the resolution is exact.  Tensoring with $M$ makes every
differential zero, whence
$\operatorname{Tor}_i^R(M,M)=M\ne0$ for every $i\geq0$.
A finite projective resolution would make these groups vanish
in sufficiently high degrees, a contradiction.  The certificate
exhausts all four ring elements and all sixty-four triples needed
for associativity and distributivity, and checks the equal image
and kernel.  This counterexample is in characteristic two and
is itself a principal ideal ring with zero divisors.
\end{proof}

\paragraph{A4.6 (Proved; constructibility and the deleted rank term).}
\label{A4.6}
Under the hypotheses of A4.2 with $R$ noetherian, each coefficient
$a_i(\mathfrak p)$ in (A4.2a) is constructible.  If $R$ is a PID
and $N(L)$ is expressed using the manuscript's multiplicities
$\mu_i$ of disks in degrees $i,i-1$, $\nu_i$ of stalks in degree
$i$, and $\beta_i=\sum_a\beta_{i,a}$ of nonunit two-term blocks
in degrees $i,i-1$, then
\[
 \operatorname{rank}_R N_i(L)
 =\mu_i+\mu_{i+1}+\beta_i+\beta_{i+1}+\nu_i.
 \tag{A4.6a}
\]
We take $\mu_0=\beta_0=0$.  The $\mu_i$ term struck from the
second rank line of 137:17 is necessary.  In the degreewise finite
PID classification, strict perfection is equivalent to finite
support of all three families $\mu_i,\beta_i,\nu_i$; derived
perfection is equivalent only to finite support of
$\beta_i,\nu_i$, with no support condition on the contractible
disk multiplicities $\mu_i$.  Thus the following sentence on
scan~17, which includes ``(str.)'' before perfection, is correct
only for strict perfection if it requires finitely many factors
of every kind.

\begin{proof}
The finite module $N_i(L)$ is finitely presented.  In a finite
presentation $R^s\xrightarrow{A}R^r\to N_i(L)\to0$, the fibre
dimension is $r-\operatorname{rank}(A\bmod\mathfrak p)$.
The locus of each matrix rank is specified by vanishing of all
minors of one size and nonvanishing of at least one minor of the
preceding size.  It is locally closed, so the dimension function
is constructible.  Only finitely many coefficients occur when
the normalization is bounded.

The PID decomposition was proved in A2.3.  A disk or nonunit
block with upper degree $j$ contributes one to each of degrees
$j$ and $j-1$, and a stalk contributes only at its own degree.
Counting these contributions proves (A4.6a).  In particular a
single disk with upper degree $i$ has rank one in degree $i$;
deleting $\mu_i$ would incorrectly give zero there.  The finite
certificate checks disks with lower degrees zero through seven
in chain degrees zero through nine.  The contribution count
proves the formula in every degree, including degree zero.
Finally, nonunit blocks contribute homology $R/(a)$ in their
lower degree, stalks contribute free homology in their own degree,
and disks contribute no homology.  This gives the two support
criteria: the strict one follows from the nonnegative rank
formula, while the derived one follows from A4.3 for the regular
noetherian ring $R$.  The zero ring is again trivial.
\end{proof}

\paragraph{Reading and literature scope.}
Our mathematical arguments use the definitions and equations of
137:14--17 and 49--56 as transcribed, and the nonnegative indexing
is unambiguous in those equations.  We do not certify an ink
reading by citing an earlier reader.  The potentially ambiguous
word ``$\Sigma$-kernel'' has been addressed in both its ordinary
and admissible senses in A2.2--A2.4.  No result above depends on
choosing an unverified alternative inside an \texttt{unsure}
mark.  We searched the cited Dold--Puppe paper, B\"uhler's
arXiv source and the Stacks Project's sections on perfect
complexes and regular rings.  We make no priority claim for the
elementary matrix formulation or the counterexamples; the
search established the cited background and did not purport to
exhaust all earlier literature on these formulations.
